\section{The OLE Formula}
\label{sec:ole}

The empirical findings converge on a three-step practical prescription. The OLE formula\footnote{Originally termed ODE (Orientation--Discipline--Expertise). Renamed to OLE after empirical field elicitation research (Papers 47--50) established that \texttt{lens} is the field name that reliably produces Tier~2 content in structured persona authoring.}
--- Orientation, Lens, Expertise --- is a minimum sufficient creative prompt
construction protocol derivable from the tier model. It specifies not what content to
include but in what order, and at what level of formulation, to construct a creative
prompt that crosses the quality threshold in any domain.

\subsection{Formalization}

The OLE formula consists of three steps, corresponding to the three tiers.

Step 1, Orientation, is a single sentence: \textit{``Think first about how the other
person will experience this.''} This sentence activates Tier 1. It names the receiver,
names the receiver's experience as the primary criterion, and establishes the evaluative
frame before any domain content enters the system. It is not a domain claim and does not
require domain knowledge to apply. The step 1 sentence is the irreplaceable component.
Every piece of evidence in this paper converges on the same finding: a creative prompt
without an orientation sentence fails at or below the bare baseline, regardless of how
much domain expertise follows it. Do not skip step 1.

Step 2, Lens, consists of two sentences: \textit{``Remove anything that is not
necessary. Verify every claim before finishing.''} These sentences activate Tier 2. The
first is an elegance gate: it enforces the operational constraint that the output contain
no unnecessary elements, which is the authoring-side analog of the reviewer's test for
over-scaffolding. The second is a rigor gate: it enforces verification of every factual
claim embedded in the output, which is the authoring-side analog of Verify the Last Mile
and the Computer Test in the reviewing profile. Together they operationalize the
orientation into checkable decision criteria. The leave-one-out data confirm their load-bearing
role: removing the elegance gate (L-T4 = $-$11.7 points) and removing the rigor gate
(L-T5 = $-$12.7 points) both drop quality below the threshold.

Step 3, Expertise, is domain content provided as context, not as persona identity. The
correct form is: \textit{``Here is the relevant domain information: [world data, style
guides, precedents, examples].''} The incorrect form is: \textit{``You are a [domain
expert].''} The distinction is not cosmetic. Providing domain content as context
activates Tier 3 in its amplifying role: the system uses the information to enrich an
output already structured by Tier 1 and Tier 2. Providing domain content as a persona
identity label activates Tier 3 first, before any orientation has been established, and
inverts the quality hierarchy. The persona study quantifies the cost of this inversion at
15.6 points on a 45-point scale.

The three-step structure is the minimum sufficient creative prompt. The reconstruction
curve from Phase 3 of the authoring study \cite{paper6} provides direct evidence: R3,
which corresponds precisely to Steps 1 and 2 without Step 3 (T1 + T4 + T5 = 35.7 out of
45), is the first condition to cross the publishable quality threshold. Every condition
without T1, regardless of what else is present, fails to cross that threshold. Every
condition with T1 and both Tier 2 traits, regardless of which Tier 3 traits are also
present, crosses the threshold. OLE Step 3 content shifts performance from 35.7 toward
42.3 --- a meaningful gain, but one that presupposes Steps 1 and 2.

\subsection{Generalization Argument}

We predict the OLE formula is not puzzle-specific (cross-domain validation pending).
The three tiers map onto structural features of quality creative work that we predict
are independent of domain.

We predict Tier 1 orientation generalizes to any AI creative task in domains where quality
is defined by receiver experience. The quality criterion of a piece of marketing copy is
whether it moves the reader toward a decision. The quality criterion of a product
specification is whether the team building from it can produce what the user needs. The
quality criterion of a legal agreement is whether the parties who sign it understand what
they have committed to. The quality criterion of training materials is whether the learner
can apply the content after instruction. In each of these cases, quality is defined by
the downstream person's experience, not by the producer's domain competence. We predict
that a marketing-expert persona activating domain conventions before audience orientation
will produce formally correct copy that fails to move readers. The tier model yields this
prediction; OLE is proposed as the correction.

We predict Tier 2 discipline generalizes similarly, because orientation must be
operationalized into checkable constraints regardless of domain. Strunk and White's
instruction to ``omit needless words'' \cite{strunkwhite} is the elegance gate expressed
for prose. Professional quality standards for factual accuracy --- from medical
record-keeping to legal cite-checking to journalistic sourcing to scientific peer review
--- are the rigor gate expressed in specific domains. We predict these constraints are not
domain-specific refinements but operational discipline that converts orientation into
defensible output; cross-domain replication is required to confirm this prediction.

Tier 3 expertise is domain-specific by definition. It is the content that makes the
creative work relevant to its particular subject matter, audience, and genre expectations.
For marketing copy, Tier 3 is product knowledge, competitor landscape, and brand voice.
For product design, it is technical specifications, component constraints, and interface
conventions. For legal drafting, it is statutory authority, precedent, and regulatory
framework. For training materials, it is the pedagogical structure of the subject and the
assumed prior knowledge of the learner. This content is necessary --- a puzzle without
world data produces puzzles disconnected from their theme; a legal brief without statutory
knowledge is merely formally structured reasoning --- but substitutable across sources
and never sufficient alone.

Table~\ref{tab:generalization} maps the OLE formula across four representative
business-domain creative tasks. The predicted tier structure follows from the receiver
criterion (Tier 1), the universal quality constraints (Tier 2), and the domain-specific
content (Tier 3) in each case.

\begin{table}[t]
\centering
\caption{Predicted OLE tier structure for representative business-domain creative tasks
(cross-domain validation pending). Tier 1 specifies the orientation sentence content;
Tier 2 specifies the discipline gates; Tier 3 specifies what domain content to provide
as context.}
\label{tab:generalization}
\begin{tabular}{p{2.0cm}p{3.2cm}p{3.2cm}p{3.2cm}}
\toprule
\textbf{Domain} & \textbf{Tier 1 (orientation)} & \textbf{Tier 2 (lens)} & \textbf{Tier 3 (expertise)} \\
\midrule
Marketing copy & How does the reader feel after reading this? & Remove non-essential claims; verify every factual assertion & Product specifications, competitor positioning, brand voice guide \\
\addlinespace
Product design & What does the user experience at each step? & Remove unnecessary interactions; verify each state is reachable & Technical specs, component constraints, interface conventions \\
\addlinespace
Legal drafting & What does the signer need to understand to comply? & Remove ambiguous phrasing; verify every obligation is enforceable & Statutory authority, controlling precedent, regulatory framework \\
\addlinespace
Training materials & What does the learner experience and retain? & Remove unnecessary content; verify every claim against the subject & Pedagogical structure, prior knowledge baseline, domain content \\
\bottomrule
\end{tabular}
\end{table}

\subsection{When Domain Expertise Is Tier 1}

The OLE formula makes a specific boundary-condition prediction: domain expertise occupies
Tier 1 only in domains where the quality criterion is defined by domain convention rather
than by receiver experience. The clearest examples are formal logic, mathematical proof,
and programming correctness. In a formal proof, ``correct'' means that each step follows
from the axioms by valid inference rules. The receiver's experience of the proof does not
define its correctness; the proof's relationship to the formal system does. Tier 1 for a
formal proof is therefore a domain criterion, not a receiver-experience criterion. OLE
does not apply to formal proof in its standard form.

For creative work broadly defined --- any work whose quality depends on what a human
receiver experiences, understands, or decides after encountering it --- we predict Tier 1
is always orientation toward the receiver. This class plausibly includes all the domains
in Table~\ref{tab:generalization} and much of the content produced in the modern economy:
creative and marketing content, interface and product design, communication and
documentation, instruction and training, policy and legal interpretation addressed to
non-specialist audiences. The boundary condition is narrow; we predict the scope of OLE's
application is wide, but this prediction requires cross-domain replication to establish
empirically.
